% Layered poisoning defenses on the repository update path (Appendix D.1).
% Single-column stack. Ink language matches Figure~\ref{fig:EmergingCaseStudy}
% and Figure~\ref{fig:EntryFormat} (parchment fill, Hold stroke, left spine).
\begin{figure}[t]
  \centering
  \resizebox{0.96\columnwidth}{!}{%
  \begin{tikzpicture}[
    font=\sffamily\scriptsize,
    card/.style={
      rounded corners=2.4pt,
      line width=0.75pt,
      draw=HoldInk,
      fill=MetaFill,
      inner xsep=5.5pt,
      inner ysep=3.4pt,
      align=left,
      font=\sffamily\scriptsize,
      text width=\dimexpr\columnwidth-40pt\relax,
    },
    idx/.style={
      circle,
      font=\sffamily\bfseries\tiny,
      inner sep=0pt,
      minimum size=10.5pt,
      text=white,
      fill=HoldInk,
    },
  ]
    \node[idx] (i1) {1};
    \node[card, right=8pt of i1] (c1) {%
      {\color{HoldInk}\ttbf{L1}}
      {\color{FigInk}\textbf{ Data Boundary}}\\[2pt]
      {\color{FigMute}Raw trigger text never enters the generator prompt.
        Only structured features and a similar-sample cluster do.}%
    };
    \node[
      font=\sffamily\scriptsize\bfseries,
      text=FigInk,
      anchor=south west,
    ] at ([yshift=5pt]c1.north west) {Local Update Path};

    \node[card, below=8pt of c1.south west, anchor=north west] (c2) {%
      {\color{HoldInk}\ttbf{L2}}
      {\color{FigInk}\textbf{ Verification Sandbox}}\\[2pt]
      {\color{FigMute}Deterministic verification is the trust anchor.
        Signatures run in a child process with a timeout and static ReDoS
        rejection.}%
    };
    \node[idx, anchor=east] (i2) at ([xshift=-8pt]c2.west) {2};

    \node[card, below=8pt of c2.south west, anchor=north west] (c3) {%
      {\color{HoldInk}\ttbf{L3}}
      {\color{FigInk}\textbf{ Reviewer Hardening}}\\[2pt]
      {\color{FigMute}The independent reviewer emits a strict JSON review
        sheet and treats every candidate as code or data, not instructions.}%
    };
    \node[idx, anchor=east] (i3) at ([xshift=-8pt]c3.west) {3};

    \node[card, below=8pt of c3.south west, anchor=north west] (c4) {%
      {\color{HoldInk}\ttbf{L4}}
      {\color{FigInk}\textbf{ Lesson Integrity}}\\[2pt]
      {\color{FigMute}Lessons are append-only, schema-validated, and
        originate only from verification or review output.}%
    };
    \node[idx, anchor=east] (i4) at ([xshift=-8pt]c4.west) {4};

    \node[card, below=8pt of c4.south west, anchor=north west] (c5) {%
      {\color{HoldInk}\ttbf{L5}}
      {\color{FigInk}\textbf{ Resource Guards}}\\[2pt]
      {\color{FigMute}Cooldown, daily evolution limits, per-run token
        budgets, and maximum rounds bound the update path.}%
    };
    \node[idx, anchor=east] (i5) at ([xshift=-8pt]c5.west) {5};

    \node[card, below=8pt of c5.south west, anchor=north west] (c6) {%
      {\color{HoldInk}\ttbf{L6}}
      {\color{FigInk}\textbf{ Retirement Protection}}\\[2pt]
      {\color{FigMute}Only negative-score or descendant-covered nodes may
        be demoted by rank, so healthy skills cannot be displaced.}%
    };
    \node[idx, anchor=east] (i6) at ([xshift=-8pt]c6.west) {6};

    \begin{scope}[on background layer]
      \draw[Spine, line width=1.2pt]
        (i1.center) -- (i2.center) -- (i3.center)
        -- (i4.center) -- (i5.center) -- (i6.center);
      % Hold west bars, matching Figure~\ref{fig:EntryFormat} cards.
      \foreach \c in {c1,c2,c3,c4,c5,c6} {
        \fill[HoldInk]
          ([xshift=0.6pt,yshift=-1.6pt]\c.north west)
          rectangle
          ([xshift=2.6pt,yshift=1.6pt]\c.south west);
      }
    \end{scope}
  \end{tikzpicture}%
  }
  \caption[Six poisoning defenses on the local update path of the defense repository.]{%
    Six poisoning defenses on the local update path of the defense
    repository.
    Layers are checked in order from L1 to L6.
  }
  \label{fig:appendix-repository-guards}
\end{figure}
